Fix \(k\) and \(h\), and abbreviate \(v=v_h\) and \(\mathcal F=\mathcal F_k^{\mathrm{tr}}\). For every \(f\in L_2(P_Z)\), the Hilbert-space square identity gives
\[
 2E_P[v(Z)f(Z)]-E_P[f(Z)^2]
 =2\langle v,f\rangle-\|f\|^2
 =\|v\|^2-\|v-f\|^2.
\]
Taking the supremum over \(f\in\mathcal F\) and subtracting from \(\|v\|^2\) yields
\[
 \|v\|^2-\sup_{f\in\mathcal F}
 \{2\langle v,f\rangle-\|f\|^2\}
 =\inf_{f\in\mathcal F}\|v-f\|^2
 =\operatorname{dist}^2(v,\mathcal F).
\]
The right-hand side is nonnegative, so the positive-part operation in the original dual definition is redundant. Taking the supremum over \(h\in\mathcal H_k\) proves the displayed identity for \(\varepsilon_{\mathrm{tr},k}^2\).

It remains to justify the exact zero characterization. Let
\[
 \mathbb B_{1,k}=\left\{u\in\mathbb R^{|J_k|}:\sum_{j\in J_k}|u_j|\le B\right\}.
\]
This set is compact. The linear map \(u\mapsto\sum_{j\in J_k}u_jg_{kj}\) is continuous from a finite-dimensional Euclidean space into \(L_2(P_Z)\), so its image \(\mathcal F_k^{\mathrm{tr}}\) is compact and hence closed. Since squared distances are nonnegative,
\[
 \sup_{h\in\mathcal H_k}
 \operatorname{dist}^2(v_h,\mathcal F_k^{\mathrm{tr}})=0
\]
holds if and only if every one of those distances is zero. Closedness then gives
\[
 \operatorname{dist}(v_h,\mathcal F_k^{\mathrm{tr}})=0
 \quad\Longleftrightarrow\quad
 v_h\in\mathcal F_k^{\mathrm{tr}}.
\]
For completeness, \(h\mapsto v_h=T(h^\dagger-h)\) is continuous by \(\texttt{lem:operator-contraction}\), and distance to a nonempty set is one-Lipschitz. Thus compactness of \(\mathcal H_k\) also makes the outer supremum a maximum. If a general critic set \(\mathcal C\) were not closed, the elementary identity \(\operatorname{dist}(v,\mathcal C)=0\iff v\in\overline{\mathcal C}\) would give only the closure statement. Finally, \(\mathcal F_k^{\mathrm{tr}}\) is generally a bounded coefficient body rather than a linear space, so membership in \(\operatorname{span}(\mathcal F_k^{\mathrm{tr}})\) alone does not force zero distance.